\documentclass[11pt]{article}
\usepackage{preamble}
\renewcommand{\answer}[1]{}

\begin{document}

\ladderheading{AP Statistics \textperiodcentered\ Unit 1 \textperiodcentered\ Topics 1.1--1.7 (CED 1.1--1.7)}{Screen Time, Two Deletions}

\focusbanner{THE PROBLEM}

A school nurse recorded, for a random sample of 40 students, three things: grade level (Freshman or Senior), the device they use most (Phone, Laptop, or Console), and daily screen time in minutes from the phone\textquotesingle s own report. The principal asked one question: \emph{Is too much screen time a concern for our students?} The nurse\textquotesingle s first report said only: \emph{12 of the 40 are over 240, that is 30\%, so yes.}\par\smallskip Group sizes: 25 freshmen, 15 seniors. Over 240 minutes: 9 freshmen, 3 seniors.\par\smallskip The two largest values in the histogram (a 900-minute and a 720-minute entry) came from two students who logged a whole weekend by mistake. The nurse deleted both and re-computed: the mean fell from 205 to 178 minutes; the median moved from 176 to 174 minutes; the headline became \emph{10 of 38 are over 240, that is 26\%.}\par
\begin{center}
\begin{tikzpicture}
\begin{axis}[width=0.96\linewidth,height=1.45in,ybar interval,x tick label as interval=false,bar shift=0pt,xmin=0,xmax=960,
 ymin=0,ymax=16,xtick={0,120,240,360,480,600,720,840,960},ytick={0,4,8,12,16},xlabel={Daily screen time (minutes)},ylabel={Students},
 tick label style={font=\scriptsize},label style={font=\scriptsize},title={All 40 students, before any deletion},title style={font=\small}]
\addplot[fill=calloutblue,draw=dayblue] coordinates {(0,3) (120,14) (240,12) (360,6) (480,3) (600,0) (720,1) (840,1) (960,0)};
\end{axis}\end{tikzpicture}
\end{center}
\begin{firsttakebox}{FIRST TAKE --- one sentence before you work the parts}
Before any calculation: does the nurse\textquotesingle s first sentence answer the principal\textquotesingle s question? Name one thing the sentence leaves out.
\workspace{0.7in}
\end{firsttakebox}

\begin{callout}{\IconBook\ RULES YOU MUST USE (NO ANSWERS HERE)}{calloutgreen}
A number means nothing until you say WHO was measured, WHAT was measured, the UNITS, and the QUESTION it answers.\\ Individuals are the things measured; variables are the characteristics recorded. A label or ID is not a measurement.\\ Groups of different sizes are compared with relative frequencies (percents), never raw counts alone.\\ A full description names shape, center, variability, AND unusual features. Some statistics resist unusual values; some do not.
\end{callout}

\newpage
\textbf{Work (a)--(d).}\par\smallskip

\dokbadge{1}~\textbf{(a)} Fill the table: who or what are the individuals, and for each recorded variable, what ONE value looks like and whether it is categorical or quantitative. Checklist: an individual is one thing measured; a name, label, or ID is not a measurement; quantitative values could sensibly be averaged.\par\smallskip \begin{tabular}{|p{1.3in}|p{2.2in}|p{1.6in}|}\hline \textbf{Individuals:} & \multicolumn{2}{l|}{\hrulefill} \\ \hline \textbf{Variable} & \textbf{One value looks like} & \textbf{Categorical / Quantitative} \\ \hline Grade level & & \\[10pt] \hline Device used most & & \\[10pt] \hline Daily screen time & & \\[10pt] \hline \end{tabular}\par
\workspace{0.15in}

\dokbadge{2}~\textbf{(b)} A senior says: \emph{Nine freshmen are over 240 and only three seniors are, so freshmen are three times as bad.} Fix the comparison. Frame: \emph{Of the \underline{\hspace{0.4in}} freshmen, \underline{\hspace{0.4in}} are over 240, which is \underline{\hspace{0.4in}}\%. Of the \underline{\hspace{0.4in}} seniors, \underline{\hspace{0.4in}} are over 240, which is \underline{\hspace{0.4in}}\%. So the fair comparison is \underline{\hspace{0.9in}} because the groups \underline{\hspace{0.9in}}.} Then explain in one sentence why the count 9 versus 3 can mislead here.\par
\workspace{0.4in}

\dokbadge{2}~\textbf{(c)} Deleting the two values moved the mean by 27 minutes and the median by 2 minutes. Which of these two statistics is resistant to unusual values? Explain what in the histogram caused the big move, using the words \emph{resistant} and \emph{unusual values}. Do not say which one the nurse should report yet.\par
\workspace{0.4in}

\dokbadge{3}~\textbf{(d)} $\star$ Judge the nurse\textquotesingle s trimmed report in three to five sentences. Use the frame below if it helps.\par
{\small\textbf{To earn an E, your answer must do ALL of these} (some = P; one or none = I):
\begin{itemize}[leftmargin=1.6em,itemsep=0pt,topsep=1pt,parsep=0pt]
  \item[{\setlength{\fboxsep}{0pt}\fbox{\rule{0pt}{0.65em}\hspace{0.65em}}}] Put the 26\% in context: who was counted, what was measured (with units), the question it answers.
  \item[{\setlength{\fboxsep}{0pt}\fbox{\rule{0pt}{0.65em}\hspace{0.65em}}}] Name the center to report and justify it with resistance, matching your answer in (c).
  \item[{\setlength{\fboxsep}{0pt}\fbox{\rule{0pt}{0.65em}\hspace{0.65em}}}] Explain why the mean moved a lot while the median barely moved.
  \item[{\setlength{\fboxsep}{0pt}\fbox{\rule{0pt}{0.65em}\hspace{0.65em}}}] Judge whether deleting the two values was responsible and what the report must SAY about them.
\end{itemize}}
\workspace{0.7in}

\begin{sentenceframebox}\raggedright\textbf{Frame:} I would report \blankt[0.8in]{statistic} because \blankt[1.5in]{reason}; removing the two values changed \blankt[0.8in]{statistic} much more than \blankt[0.8in]{statistic}, so the principal should \blankt[1.8in]{action}.\end{sentenceframebox}

\begin{wordbankbox}\raggedright
\textbf{Word bank:} median \ensuremath{\;\cdot\;} it resists unusual values \ensuremath{\;\cdot\;} mode \ensuremath{\;\cdot\;} mean \ensuremath{\;\cdot\;} it follows unusual values \ensuremath{\;\cdot\;} omit the deletion details \ensuremath{\;\cdot\;} disclose and justify the deletions
\par\smallskip{\footnotesize\color{framegray} Some words are not needed.}
\end{wordbankbox}

\textbf{Turn this sheet in whenever you finish --- bonus credit. Part (d) is scored E / P / I.}\par

\end{document}
